\subsection{Contexte}

Les ondes sont des perturbations qui se propagent dans l'espace en transportant de l'énergie, sans transport de matière. On distingue notamment les ondes mécaniques, comme le son, qui ont besoin d'un milieu matériel pour se propager, et les ondes électromagnétiques, comme les micro-ondes, qui se propagent aussi dans le vide \citeit{openstax_sound}. Malgré leur nature différente, ces ondes partagent des comportements communs : réflexion, réfraction, diffraction et interférences, ainsi qu'une même relation entre leur vitesse, leur longueur d'onde et leur fréquence.

Dans ce laboratoire, deux manipulations ont été réalisées afin d'étudier ces phénomènes. Dans un premier temps, un émetteur de micro-ondes de 9,4~GHz a été utilisé pour observer la propagation rectiligne, l'absorption, la réflexion, la réfraction, la diffraction, les interférences et la polarisation de ces ondes, ainsi que pour déterminer leur longueur d'onde. Dans un second temps, le trombone de Koenig a permis de mesurer la vitesse du son dans différents gaz, en exploitant les interférences entre deux ondes sonores de fréquence connue.

L'ensemble des résultats obtenus a ensuite été comparé aux modèles théoriques calculés afin d'en évaluer la validité et de mettre en évidence les sources d'erreurs expérimentales.

La réalisation de ce laboratoire s'est déroulée en salle \salle, sur une durée de 8 périodes.
